% Options for packages loaded elsewhere
% Options for packages loaded elsewhere
\PassOptionsToPackage{unicode}{hyperref}
\PassOptionsToPackage{hyphens}{url}
\PassOptionsToPackage{dvipsnames,svgnames,x11names}{xcolor}
%
\documentclass[
  polish,
  10pt,
  a4paper,
  DIV=11,
  numbers=noendperiod]{scrartcl}
\usepackage{xcolor}
\usepackage[margin=1.6cm, top=1.6cm, bottom=1.8cm]{geometry}
\usepackage{amsmath,amssymb}
\setcounter{secnumdepth}{-\maxdimen} % remove section numbering
\usepackage{iftex}
\ifPDFTeX
  \usepackage[T1]{fontenc}
  \usepackage[utf8]{inputenc}
  \usepackage{textcomp} % provide euro and other symbols
\else % if luatex or xetex
  \usepackage{unicode-math} % this also loads fontspec
  \defaultfontfeatures{Scale=MatchLowercase}
  \defaultfontfeatures[\rmfamily]{Ligatures=TeX,Scale=1}
\fi
\usepackage{lmodern}
\ifPDFTeX\else
  % xetex/luatex font selection
  \setmainfont[]{Lato}
\fi
% Use upquote if available, for straight quotes in verbatim environments
\IfFileExists{upquote.sty}{\usepackage{upquote}}{}
\IfFileExists{microtype.sty}{% use microtype if available
  \usepackage[]{microtype}
  \UseMicrotypeSet[protrusion]{basicmath} % disable protrusion for tt fonts
}{}
\makeatletter
\@ifundefined{KOMAClassName}{% if non-KOMA class
  \IfFileExists{parskip.sty}{%
    \usepackage{parskip}
  }{% else
    \setlength{\parindent}{0pt}
    \setlength{\parskip}{6pt plus 2pt minus 1pt}}
}{% if KOMA class
  \KOMAoptions{parskip=half}}
\makeatother
% Make \paragraph and \subparagraph free-standing
\makeatletter
\ifx\paragraph\undefined\else
  \let\oldparagraph\paragraph
  \renewcommand{\paragraph}{
    \@ifstar
      \xxxParagraphStar
      \xxxParagraphNoStar
  }
  \newcommand{\xxxParagraphStar}[1]{\oldparagraph*{#1}\mbox{}}
  \newcommand{\xxxParagraphNoStar}[1]{\oldparagraph{#1}\mbox{}}
\fi
\ifx\subparagraph\undefined\else
  \let\oldsubparagraph\subparagraph
  \renewcommand{\subparagraph}{
    \@ifstar
      \xxxSubParagraphStar
      \xxxSubParagraphNoStar
  }
  \newcommand{\xxxSubParagraphStar}[1]{\oldsubparagraph*{#1}\mbox{}}
  \newcommand{\xxxSubParagraphNoStar}[1]{\oldsubparagraph{#1}\mbox{}}
\fi
\makeatother


\usepackage{longtable,booktabs,array}
\newcounter{none} % for unnumbered tables
\usepackage{calc} % for calculating minipage widths
% Correct order of tables after \paragraph or \subparagraph
\usepackage{etoolbox}
\makeatletter
\patchcmd\longtable{\par}{\if@noskipsec\mbox{}\fi\par}{}{}
\makeatother
% Allow footnotes in longtable head/foot
\IfFileExists{footnotehyper.sty}{\usepackage{footnotehyper}}{\usepackage{footnote}}
\makesavenoteenv{longtable}
\usepackage{graphicx}
\makeatletter
\newsavebox\pandoc@box
\newcommand*\pandocbounded[1]{% scales image to fit in text height/width
  \sbox\pandoc@box{#1}%
  \Gscale@div\@tempa{\textheight}{\dimexpr\ht\pandoc@box+\dp\pandoc@box\relax}%
  \Gscale@div\@tempb{\linewidth}{\wd\pandoc@box}%
  \ifdim\@tempb\p@<\@tempa\p@\let\@tempa\@tempb\fi% select the smaller of both
  \ifdim\@tempa\p@<\p@\scalebox{\@tempa}{\usebox\pandoc@box}%
  \else\usebox{\pandoc@box}%
  \fi%
}
% Set default figure placement to htbp
\def\fps@figure{htbp}
\makeatother



\ifLuaTeX
\usepackage[bidi=basic,shorthands=off]{babel}
\else
\usepackage[bidi=default,shorthands=off]{babel}
\fi
\ifPDFTeX
\else
\babelfont{rm}[]{Lato}
\fi
\ifLuaTeX
  \usepackage{selnolig} % disable illegal ligatures
\fi


\setlength{\emergencystretch}{3em} % prevent overfull lines

\providecommand{\tightlist}{%
  \setlength{\itemsep}{0pt}\setlength{\parskip}{0pt}}



 


\usepackage{xcolor}
\definecolor{upwraccent}{HTML}{6b1a2a}
\definecolor{upwrpanel}{HTML}{faf6ec}
\definecolor{upwrrule}{HTML}{d9cfbc}
\definecolor{upwrsoft}{HTML}{3d3832}
\usepackage[most]{tcolorbox}
\usepackage{enumitem}
\usepackage{longtable,booktabs,array}
\usepackage{needspace}
\newfontfamily\symfont{DejaVu Sans}
\setlist{nosep, leftmargin=1.4em}
\setlength{\parindent}{0pt}
\setkomafont{section}{\Large\bfseries\color{upwraccent}}
\setkomafont{subsection}{\large\bfseries\color{upwraccent}}
\newtcolorbox{jbox}[1]{enhanced, breakable=false, colback=upwrpanel,
  colframe=upwrrule, boxrule=0.4pt, arc=1.5pt, left=5pt, right=5pt,
  top=3pt, bottom=3pt, before skip=6pt, after skip=6pt,
  title={#1}, coltitle=upwraccent, colbacktitle=upwrpanel,
  fonttitle=\bfseries, titlerule=0pt, toptitle=2pt, bottomtitle=0pt}
\newcommand{\jmenu}[1]{{\color{upwraccent}\small #1}}
\KOMAoption{captions}{tableheading}
\makeatletter
\@ifpackageloaded{caption}{}{\usepackage{caption}}
\AtBeginDocument{%
\ifdefined\contentsname
  \renewcommand*\contentsname{Spis treści}
\else
  \newcommand\contentsname{Spis treści}
\fi
\ifdefined\listfigurename
  \renewcommand*\listfigurename{Spis rycin}
\else
  \newcommand\listfigurename{Spis rycin}
\fi
\ifdefined\listtablename
  \renewcommand*\listtablename{Spis tabel}
\else
  \newcommand\listtablename{Spis tabel}
\fi
\ifdefined\figurename
  \renewcommand*\figurename{Rysunek}
\else
  \newcommand\figurename{Rysunek}
\fi
\ifdefined\tablename
  \renewcommand*\tablename{Tabela}
\else
  \newcommand\tablename{Tabela}
\fi
}
\@ifpackageloaded{float}{}{\usepackage{float}}
\floatstyle{ruled}
\@ifundefined{c@chapter}{\newfloat{codelisting}{h}{lop}}{\newfloat{codelisting}{h}{lop}[chapter]}
\floatname{codelisting}{Wykaz}
\newcommand*\listoflistings{\listof{codelisting}{Spis wykazów}}
\makeatother
\makeatletter
\makeatother
\makeatletter
\@ifpackageloaded{caption}{}{\usepackage{caption}}
\@ifpackageloaded{subcaption}{}{\usepackage{subcaption}}
\makeatother
\usepackage{bookmark}
\IfFileExists{xurl.sty}{\usepackage{xurl}}{} % add URL line breaks if available
\urlstyle{same}
% fallback for those not using the hyperref driver hyperxmp:
\makeatletter
\@ifundefined{xmpquote}{\newcommand{\xmpquote}[1]{#1}}{}
\makeatother
\hypersetup{
  pdftitle={Ściąga jUPWR},
  pdflang={pl},
  colorlinks=true,
  linkcolor={upwraccent},
  filecolor={Maroon},
  citecolor={Blue},
  urlcolor={Blue},
  pdfcreator={LaTeX via pandoc}}


\title{Ściąga jUPWR}
\usepackage{etoolbox}
\makeatletter
\providecommand{\subtitle}[1]{% add subtitle to \maketitle
  \apptocmd{\@title}{\par {\large #1 \par}}{}{}
}
\makeatother
\subtitle{Gdzie w jUPWR robi się to, czego potrzeba na wykładzie ze
statystyki}
\author{}
\date{}
\begin{document}
\maketitle


\textit{Stan na jUPWR 1.0.5. Wpis każdego tematu jest raz — przy pierwszym wykładzie, na którym się pojawia; dalsze wykłady odsyłają do strony.}

\subsection{Mapa menu}\label{mapa-menu}

\begin{jbox}{Wstążka jUPWR}
\textbf{Zakładki:} Zmienne · Dane · Analizy · Wykresy · Edytuj. \quad Plik: przycisk {\symfont ☰} w lewym górnym rogu.

\smallskip
\begin{tabular}{@{}>{\bfseries\color{upwraccent}}p{3.6cm}p{13cm}@{}}
Eksploracja & Zmienne ilościowe · Zmienne jakościowe · Szereg rozdzielczy \\
Rozkłady & rozkłady ciągłe (normalny, t, χ², F, wykładniczy …) i dyskretne (dwumianowy, Poissona …) \\
Przedziały ufności & Średnie · Proporcje · Korelacja i regresja · Dydaktyka (Jak działa bootstrap) \\
Testy t & Jedna próba · Dwie grupy · Próby sparowane \\
ANOVA & ANOVA · ANOVA powtórzonych pomiarów \\
Regresja & Korelacja · Korelacja cząstkowa · Regresja liniowa · Regresja logistyczna \\
Częstości & Jedna zmienna · Tabela kontyngencji · Próby zależne \\
Testy permutacyjne & Jedna próba · Dwie grupy · Próby sparowane \\
\end{tabular}

\smallskip
\textbf{Panel analizy:} najważniejsze opcje są widoczne od razu; reszta w zwiniętych sekcjach \textbf{Założenia} i \textbf{Zaawansowane}. Na dole każdej analizy: \textbf{Opis zastosowanych metod} (do raportu).
\end{jbox}

\section{Część A --- wg
wykładów}\label{czux119ux15bux107-a-wg-wykux142aduxf3w}

\needspace{8\baselineskip}

\subsection{Wykład 01 --- Statystyka
opisowa}\label{wykux142ad-01-statystyka-opisowa}

\begin{jbox}{Otwarcie pliku z danymi}\label{w:otworz-plik}
\jmenu{\textbf{{\symfont ☰} (menu pliku)} → \textbf{Otwórz} → \textbf{Przeglądaj}}

\begin{enumerate}
\item Wybierz plik CSV, XLSX albo zapisany projekt .omv.
\item Gotowe zbiory z zajęć: {\symfont ☰} → Otwórz → Biblioteka.
\end{enumerate}
\textit{Wynik:} Dane pojawiają się w arkuszu na zakładce Dane.

\end{jbox}

\begin{jbox}{Typ zmiennej (poziom pomiaru)}\label{w:poziom-pomiaru}
\jmenu{\textbf{Dane} → \textbf{Edytuj}}

\begin{enumerate}
\item Kliknij nagłówek kolumny, potem Edytuj (albo kliknij nagłówek dwukrotnie).
\item Poziom pomiaru: Nominalna (kategorie bez porządku), Porządkowa (kategorie z porządkiem), Ciągła (liczby).
\item Typ danych: Liczba całkowita / Dziesiętny / Tekst.
\end{enumerate}
\textit{Wynik:} Ikona w nagłówku kolumny pokazuje poziom pomiaru; analizy przyjmują tylko zmienne właściwego typu.

{\color{upwrsoft}\small\textit{Uwaga:} Liczby zapisane jako tekst (np. „3-4h”, przecinek dziesiętny) trafiają jako Nominalna — popraw dane, zanim zmienisz typ.}
\end{jbox}

\begin{jbox}{Tabela liczności zmiennej jakościowej}\label{w:tabela-licznosci}
\jmenu{\textbf{Analizy} → \textbf{Eksploracja} → \textbf{Zmienne jakościowe}}

\begin{itemize}
\item[$\triangleright$] \textbf{Zmienne}: zmienna jakościowa
\end{itemize}
\begin{enumerate}
\item Tabela: Liczności i procenty skumulowane.
\item Wykres: Słupkowy.
\end{enumerate}
\textit{Wynik:} Tabela liczności z procentami; dominanta i liczba kategorii w tabeli Podsumowanie.

{\color{upwrsoft}\small\textit{Uwaga:} Wykresu kołowego w jUPWR nie ma celowo — słupkowy jest czytelniejszy.}
\end{jbox}

\begin{jbox}{Tabela krzyżowa (opis dwóch zmiennych jakościowych)}\label{w:tabela-krzyzowa-opis}
\jmenu{\textbf{Analizy} → \textbf{Eksploracja} → \textbf{Zmienne jakościowe}}

\begin{itemize}
\item[$\triangleright$] \textbf{Zmienne}: pierwsza zmienna jakościowa
\item[$\triangleright$] \textbf{Podziel według}: druga zmienna jakościowa
\end{itemize}
\begin{enumerate}
\item Procenty (przy podziale): \% w kolumnie, \% w wierszu albo \% ogółem.
\item Wykres: Słupkowy albo Mozaikowy.
\end{enumerate}
\textit{Wynik:} Tabela krzyżowa liczności z wybranymi procentami.

{\color{upwrsoft}\small\textit{Uwaga:} Test χ² dla tej tabeli jest w Częstości → Tabela kontyngencji.}
\end{jbox}

\begin{jbox}{Średnia, mediana, kwartyle, odchylenie standardowe}\label{w:statystyki-opisowe}
\jmenu{\textbf{Analizy} → \textbf{Eksploracja} → \textbf{Zmienne ilościowe}}

\begin{itemize}
\item[$\triangleright$] \textbf{Zmienne}: jedna lub kilka zmiennych ilościowych
\end{itemize}
\begin{enumerate}
\item Domyślnie: N, Braki, Średnia, Mediana, Kwartyle, Odchylenie standardowe, Minimum, Maksimum, Współczynnik zmienności V (\%).
\item Zaawansowane → Rozproszenie: Wariancja, Rozstęp, Rozstęp ćwiartkowy (IQR).
\item Zaawansowane → Położenie: Dominanta. Zaawansowane → Kwantyle: Percentyle (np. 10, 90).
\end{enumerate}
\textit{Wynik:} Tabela statystyk; Układ tabeli przełącza zmienne w wierszach / w kolumnach.

\end{jbox}

\begin{jbox}{Statystyki opisowe osobno w grupach}\label{w:porownanie-grup-opis}
\jmenu{\textbf{Analizy} → \textbf{Eksploracja} → \textbf{Zmienne ilościowe}}

\begin{itemize}
\item[$\triangleright$] \textbf{Zmienne}: zmienna ilościowa
\item[$\triangleright$] \textbf{Podziel według}: zmienna grupująca (można kilka)
\end{itemize}
\begin{enumerate}
\item Wykres: Pudełkowy — pudełka grup obok siebie.
\end{enumerate}
\textit{Wynik:} Jeden wiersz (lub kolumna) statystyk na grupę.

\end{jbox}

\begin{jbox}{Histogram}\label{w:histogram}
\jmenu{\textbf{Analizy} → \textbf{Eksploracja} → \textbf{Zmienne ilościowe}}

\begin{itemize}
\item[$\triangleright$] \textbf{Zmienne}: zmienna ilościowa
\end{itemize}
\begin{enumerate}
\item Wykres: Histogram.
\item Zaawansowane → Wykresy dodatkowe: Gęstość na histogramie.
\end{enumerate}
\textit{Wynik:} Histogram pod tabelą statystyk.

{\color{upwrsoft}\small\textit{Uwaga:} Osobny, bardziej konfigurowalny histogram: Wykresy → Rozkład → Histogram.}
\end{jbox}

\begin{jbox}{Wykres pudełkowy}\label{w:pudelkowy}
\jmenu{\textbf{Analizy} → \textbf{Eksploracja} → \textbf{Zmienne ilościowe}}

\begin{itemize}
\item[$\triangleright$] \textbf{Zmienne}: zmienna ilościowa
\item[$\triangleright$] \textbf{Podziel według}: (opcjonalnie) grupy
\end{itemize}
\begin{enumerate}
\item Wykres: Pudełkowy (Zaznacz średnią jest włączone).
\item Zaawansowane → Wykresy dodatkowe: Punkty na pudełkowym, Skrzypce na pudełkowym.
\end{enumerate}
\textit{Wynik:} Pudełko = Q1–Q3, kreska = mediana, punkty poza wąsami = wartości odstające.

\end{jbox}

\begin{jbox}{Skośność i kurtoza}\label{w:ksztalt}
\jmenu{\textbf{Analizy} → \textbf{Eksploracja} → \textbf{Zmienne ilościowe}}

\begin{itemize}
\item[$\triangleright$] \textbf{Zmienne}: zmienna ilościowa
\end{itemize}
\begin{enumerate}
\item Zaawansowane → Kształt: Skośność, Kurtoza.
\end{enumerate}
\textit{Wynik:} Kolumny/wiersze Skośność i Kurtoza w tabeli statystyk (z błędem standardowym).

\end{jbox}

\begin{jbox}{Szereg rozdzielczy (dane w klasach)}\label{w:szereg-rozdzielczy}
\jmenu{\textbf{Analizy} → \textbf{Eksploracja} → \textbf{Szereg rozdzielczy}}

\begin{itemize}
\item[$\triangleright$] \textbf{Zmienna}: zmienna ilościowa
\end{itemize}
\begin{enumerate}
\item Podział na klasy: Reguła Sturgesa, Liczba klas albo Szerokość klasy.
\item Wykres: Histogram klas, Krzywa częstości skumulowanych.
\end{enumerate}
\textit{Wynik:} Tabela klas z licznościami i częstościami skumulowanymi; miary z szeregu obok wartości z danych.

\end{jbox}

\needspace{8\baselineskip}

\subsection{Wykład 02 --- Rozkłady
prawdopodobieństwa}\label{wykux142ad-02-rozkux142ady-prawdopodobieux144stwa}

\begin{jbox}{Prawdopodobieństwo i kwantyl w rozkładzie normalnym}\label{w:rozklad-normalny}
\jmenu{\textbf{Analizy} → \textbf{Rozkłady} → \textbf{Rozkład normalny}}

\begin{enumerate}
\item Parametry: Średnia =, Odch. std. = (dla wartości z: 0 i 1).
\item Oblicz prawdopodobieństwo: P(X ≤ x1), P(X ≥ x1) albo P(x1 ≤ X ≤ x2); wpisz x1 (i x2).
\item Oblicz kwantyl(e): kwantyl kumulatywny albo kwantyle przedziału centralnego dla p =
\end{enumerate}
\textit{Wynik:} Tabela z prawdopodobieństwem / kwantylem i wykres gęstości z zakreślonym obszarem.

\end{jbox}

\begin{jbox}{Rozkład dwumianowy, Poissona, geometryczny}\label{w:rozklady-dyskretne}
\jmenu{\textbf{Analizy} → \textbf{Rozkłady} → \textbf{Rozkład dwumianowy / Poissona / geometryczny}}

\begin{enumerate}
\item Parametry rozkładu (dwumianowy: n =, Prawdopodb. =).
\item Oblicz prawdopodobieństwo: P(X = x1), P(X ≤ x1), P(X ≥ x1) albo P(x1 ≤ X ≤ x2).
\item Momenty: Pokaż wartość oczekiwaną E[X], Pokaż wariancję Var[X].
\end{enumerate}
\textit{Wynik:} Tabela z prawdopodobieństwem i wykres funkcji prawdopodobieństwa.

\end{jbox}

\begin{jbox}{Rozkłady t, χ², F, wykładniczy}\label{w:rozklady-ciagle}
\jmenu{\textbf{Analizy} → \textbf{Rozkłady} → \textbf{Rozkład t / χ² / F / wykładniczy}}

\begin{enumerate}
\item Ten sam układ co rozkład normalny: parametry → Oblicz prawdopodobieństwo albo Oblicz kwantyl(e).
\item Przykład: wartość krytyczna t dla 95\% — rozkład t, kwantyle przedziału centralnego, p = 0.95.
\end{enumerate}
\textit{Wynik:} Tabela z prawdopodobieństwem / kwantylem i wykres gęstości.

\end{jbox}

\needspace{8\baselineskip}

\subsection{Wykład 03 --- Przedziały
ufności}\label{wykux142ad-03-przedziaux142y-ufnoux15bci}

\begin{jbox}{Przedział ufności dla średniej}\label{w:ci-srednia}
\jmenu{\textbf{Analizy} → \textbf{Przedziały ufności} → \textbf{Średnie} → \textbf{Średnia}}

\begin{itemize}
\item[$\triangleright$] \textbf{Zmienna}: zmienna ilościowa
\item[$\triangleright$] \textbf{Grupy (opcjonalnie)}: osobny przedział w każdej grupie
\end{itemize}
\begin{enumerate}
\item Poziom ufności = 95 (zmień, żeby zobaczyć wpływ na szerokość).
\item Metoda przedziału: t-Studenta.
\end{enumerate}
\textit{Wynik:} Tabela: średnia i granice przedziału; wykres przedziału.

\end{jbox}

\begin{jbox}{Przedział ufności dla różnicy średnich (dwie grupy)}\label{w:ci-roznica-srednich}
\jmenu{\textbf{Analizy} → \textbf{Przedziały ufności} → \textbf{Średnie} → \textbf{Różnica średnich}}

\begin{itemize}
\item[$\triangleright$] \textbf{Zmienna}: zmienna ilościowa
\item[$\triangleright$] \textbf{Zmienna grupująca}: dwie grupy (przy więcej: wybierz Grupa 1 i Grupa 2)
\end{itemize}
\begin{enumerate}
\item Metoda przedziału: Welcha (nierówne wariancje) — domyślnie; t-Studenta (wspólna wariancja) dla wersji klasycznej.
\item Wielkość efektu: d Cohena z przedziałem ufności.
\end{enumerate}
\textit{Wynik:} Tabela: różnica średnich z przedziałem (i d Cohena z przedziałem); wykres różnicy.

\end{jbox}

\begin{jbox}{Przedział ufności dla średniej różnic (pomiary sparowane)}\label{w:ci-srednia-roznic}
\jmenu{\textbf{Analizy} → \textbf{Przedziały ufności} → \textbf{Średnie} → \textbf{Średnia różnic}}

\begin{itemize}
\item[$\triangleright$] \textbf{Pomiar 1}: pierwszy pomiar (kolumna)
\item[$\triangleright$] \textbf{Pomiar 2}: drugi pomiar (kolumna)
\end{itemize}
\begin{enumerate}
\item Metoda przedziału: t-Studenta. Opcjonalnie: d Cohena z przedziałem ufności.
\end{enumerate}
\textit{Wynik:} Tabela: średnia różnic z przedziałem.

\end{jbox}

\begin{jbox}{Przedział ufności dla proporcji}\label{w:ci-proporcja}
\jmenu{\textbf{Analizy} → \textbf{Przedziały ufności} → \textbf{Proporcje} → \textbf{Proporcja}}

\begin{itemize}
\item[$\triangleright$] \textbf{Zmienna}: zmienna jakościowa
\end{itemize}
\begin{enumerate}
\item Kategoria „sukces”: wybierz kategorię, której udział liczysz.
\item Metoda przedziału: Wilsona (score) — domyślnie; Cloppera-Pearsona (dokładna); Walda (wzór z wykładu).
\end{enumerate}
\textit{Wynik:} Tabela: proporcja z próby i granice przedziału; wykres.

{\color{upwrsoft}\small\textit{Uwaga:} Przedział Walda przy małym n lub proporcji blisko 0 albo 1 bywa za wąski albo wychodzi poza [0, 1].}
\end{jbox}

\begin{jbox}{Przedział ufności dla różnicy proporcji}\label{w:ci-roznica-proporcji}
\jmenu{\textbf{Analizy} → \textbf{Przedziały ufności} → \textbf{Proporcje} → \textbf{Różnica proporcji}}

\begin{itemize}
\item[$\triangleright$] \textbf{Zmienna}: zmienna jakościowa (wynik)
\item[$\triangleright$] \textbf{Zmienna grupująca}: dwie grupy
\end{itemize}
\begin{enumerate}
\item Kategoria „sukces”, Grupa 1, Grupa 2.
\item Metoda przedziału: Newcombe’a (score) — domyślnie; Walda.
\end{enumerate}
\textit{Wynik:} Tabela: różnica proporcji z przedziałem; wykres.

\end{jbox}

\begin{jbox}{Przedział bootstrapowy}\label{w:bootstrap}
\jmenu{\textbf{Analizy} → \textbf{Przedziały ufności} → \textbf{(dowolna analiza)}}

\begin{enumerate}
\item Metoda przedziału: Bootstrap percentylowy albo Bootstrap BCa.
\item Zaawansowane → Bootstrap: Liczba losowań B, Ziarno generatora (powtarzalny wynik), Histogram rozkładu bootstrapowego.
\item Idea krok po kroku: Przedziały ufności → Dydaktyka → Jak działa bootstrap.
\end{enumerate}
\textit{Wynik:} Ta sama tabela co dla metody klasycznej.

\end{jbox}

\needspace{8\baselineskip}

\subsection{Wykład 04 --- Testowanie
hipotez}\label{wykux142ad-04-testowanie-hipotez}

\begin{jbox}{Nowa zmienna z formuły}\label{w:oblicz-zmienna}
\jmenu{\textbf{Dane} → \textbf{Oblicz}}

\begin{enumerate}
\item Nadaj nazwę nowej zmiennej.
\item Wpisz formułę, np. english \textgreater{} 20, (read + math) / 2, Z(wzrost), read \textgreater{}= VMED(read).
\item Przycisk fx pokazuje listę funkcji.
\end{enumerate}
\textit{Wynik:} Nowa kolumna na końcu arkusza (lub obok zaznaczonej), przelicza się sama po zmianie danych.

\end{jbox}

\begin{jbox}{Test t dla jednej próby}\label{w:t-jedna-proba}
\jmenu{\textbf{Analizy} → \textbf{Testy t} → \textbf{Jedna próba}}

\begin{itemize}
\item[$\triangleright$] \textbf{Zmienne}: zmienna ilościowa
\end{itemize}
\begin{enumerate}
\item Hipoteza: Wartość testowa = (np. 650).
\item Hipoteza alternatywna: Dwustronna (≠), Większa (\textgreater{}) albo Mniejsza (\textless{}).
\item Wykres pudełkowy z punktami jest włączony; linia pokazuje wartość testową.
\end{enumerate}
\textit{Wynik:} Tabela testu: Statystyka t, df, p, różnica średnich, wielkość efektu (d Cohena).

{\color{upwrsoft}\small\textit{Uwaga:} Przedział ufności dla średniej: Przedziały ufności → Średnie → Średnia.}
\end{jbox}

\begin{jbox}{Test dla jednej proporcji (dwumianowy)}\label{w:test-proporcji}
\jmenu{\textbf{Analizy} → \textbf{Częstości} → \textbf{Jedna zmienna}}

\begin{itemize}
\item[$\triangleright$] \textbf{Zmienna}: zmienna o dwóch kategoriach
\end{itemize}
\begin{enumerate}
\item Proporcje oczekiwane: wpisz proporcję z H₀ (np. 0,5 i 0,5).
\item Przy dwóch kategoriach test dwumianowy (dokładny) wybiera się sam.
\item Zaawansowane → Hipoteza (2 kategorie): Hipoteza alternatywna.
\end{enumerate}
\textit{Wynik:} Tabela testu z p; Wykres: Obserwowane i oczekiwane.

{\color{upwrsoft}\small\textit{Uwaga:} Przedział ufności dla proporcji: Przedziały ufności → Proporcje → Proporcja (Cloppera-Pearsona = dokładny).}
\end{jbox}

\begin{jbox}{Test χ² zgodności}\label{w:chi2-zgodnosc}
\jmenu{\textbf{Analizy} → \textbf{Częstości} → \textbf{Jedna zmienna}}

\begin{itemize}
\item[$\triangleright$] \textbf{Zmienna}: zmienna o trzech lub więcej kategoriach
\end{itemize}
\begin{enumerate}
\item Proporcje oczekiwane: wpisz udziały z H₀ (domyślnie równe).
\item Wielkość efektu (w Cohena) — włączona.
\item Zaawansowane → Post-hoc: Reszty standaryzowane (które kategorie odbiegają).
\end{enumerate}
\textit{Wynik:} Tabela testu χ² (χ², df, p) i w Cohena.

\end{jbox}

\begin{jbox}{Test χ² niezależności (tabela kontyngencji)}\label{w:chi2-niezaleznosc}
\jmenu{\textbf{Analizy} → \textbf{Częstości} → \textbf{Tabela kontyngencji}}

\begin{itemize}
\item[$\triangleright$] \textbf{Wiersze}: pierwsza zmienna jakościowa
\item[$\triangleright$] \textbf{Kolumny}: druga zmienna jakościowa
\end{itemize}
\begin{enumerate}
\item Test i wielkość efektu: χ² niezależności i Wielkość efektu (V Craméra) — włączone.
\item Procenty: Wierszami (domyślnie) albo Kolumnami — w kierunku, w którym porównujesz grupy.
\item Zaawansowane → Post-hoc: Reszty standaryzowane (które komórki decydują).
\item Wykres: Słupkowy albo Mozaikowy.
\end{enumerate}
\textit{Wynik:} Tabela liczności z procentami, tabela testu (χ², df, p) i V Craméra.

\end{jbox}

\begin{jbox}{Mała tabela — liczebności oczekiwane i test Fishera}\label{w:fisher}
\jmenu{\textbf{Analizy} → \textbf{Częstości} → \textbf{Tabela kontyngencji}}

\begin{itemize}
\item[$\triangleright$] \textbf{Wiersze}: zmienna jakościowa
\item[$\triangleright$] \textbf{Kolumny}: zmienna jakościowa
\end{itemize}
\begin{enumerate}
\item Zaawansowane → Liczności: Liczebności oczekiwane.
\item Test i wielkość efektu: Dokładny test Fishera.
\item Zaawansowane → Dodatkowe testy: Poprawka ciągłości (2×2) — poprawka Yatesa.
\end{enumerate}
\textit{Wynik:} Gdy liczebności oczekiwane są za małe (E \textless{} 5), pod tabelą testu pojawia się ostrzeżenie; wynik Fishera jest w tabeli testów.

\end{jbox}

\begin{jbox}{Test t dla dwóch grup niezależnych}\label{w:t-dwie-grupy}
\jmenu{\textbf{Analizy} → \textbf{Testy t} → \textbf{Dwie grupy}}

\begin{itemize}
\item[$\triangleright$] \textbf{Zmienne zależne}: zmienna ilościowa
\item[$\triangleright$] \textbf{Zmienna grupująca}: zmienna o dwóch kategoriach
\end{itemize}
\begin{enumerate}
\item Testy: t Welcha (nierówne wariancje); t Studenta jest zaznaczony domyślnie — odznacz, jeśli niepotrzebny.
\item Hipoteza alternatywna: Dwustronna (≠), Większa (\textgreater{}) albo Mniejsza (\textless{}).
\item Sprawdzenie założeń: Jednorodność wariancji (Levene), Normalność (Shapiro-Wilk).
\end{enumerate}
\textit{Wynik:} Tabela testu: Statystyka t, df, p, różnica średnich, wielkość efektu (d Cohena) — po jednym wierszu na test.

{\color{upwrsoft}\small\textit{Uwaga:} Przedział ufności dla różnicy średnich i d Cohena z przedziałem: Przedziały ufności → Średnie → Różnica średnich.}
\end{jbox}

\begin{jbox}{Test t dla prób sparowanych}\label{w:t-sparowane}
\jmenu{\textbf{Analizy} → \textbf{Testy t} → \textbf{Próby sparowane}}

\begin{itemize}
\item[$\triangleright$] \textbf{Pary zmiennych}: dwie kolumny: pomiar 1 i pomiar 2
\end{itemize}
\begin{enumerate}
\item Testy: t Studenta (domyślnie); Wilcoxona (rangowanych znaków) jako wersja nieparametryczna.
\item Sprawdzenie założeń: Normalność (Shapiro-Wilk) — dotyczy różnic.
\end{enumerate}
\textit{Wynik:} Tabela testu: t, df, p, średnia różnica, d Cohena; wykres z trzecim pudełkiem różnic.

{\color{upwrsoft}\small\textit{Uwaga:} Dane muszą być w formacie szerokim: jeden wiersz = jedna osoba, dwa pomiary w dwóch kolumnach.}
\end{jbox}

\begin{jbox}{ANOVA jednoczynnikowa i porównania post-hoc}\label{w:anova}
\jmenu{\textbf{Analizy} → \textbf{ANOVA} → \textbf{ANOVA}}

\begin{itemize}
\item[$\triangleright$] \textbf{Zmienna zależna}: zmienna ilościowa
\item[$\triangleright$] \textbf{Czynniki}: zmienna grupująca (trzy lub więcej grup)
\end{itemize}
\begin{enumerate}
\item Test: Wielkość efektu (η²p) — przy jednym czynniku równa η².
\item Test: Nierówne wariancje (Welch), gdy wariancje w grupach się różnią.
\item Porównania: metoda Tukeya (HSD) domyślnie; Tabela porównań parami.
\item Założenia: Jednorodność wariancji, Normalność reszt, Wykres Q-Q reszt.
\end{enumerate}
\textit{Wynik:} Tabela ANOVA (F, df, p, η²p); wykres średnich z literami — grupy z tą samą literą nie różnią się istotnie.

\end{jbox}

\begin{jbox}{Gdzie są wielkości efektu}\label{w:sila-efektu}
\jmenu{\textbf{Analizy} → \textbf{(zależnie od testu)}}

\begin{enumerate}
\item d Cohena: Testy t (w tabeli testu); z przedziałem: Przedziały ufności → Różnica średnich / Średnia różnic.
\item r: Regresja → Korelacja (sam współczynnik jest wielkością efektu).
\item V Craméra: Częstości → Tabela kontyngencji (włączone). w Cohena: Częstości → Jedna zmienna.
\item η²: ANOVA → Wielkość efektu (η²p); Zaawansowane → Model: η², ω².
\end{enumerate}
\textit{Wynik:} Kolumna wielkości efektu w tabeli danego testu.

\end{jbox}

\begin{jbox}{Korelacja (Pearson, Spearman) i macierz korelacji}\label{w:korelacja}
\jmenu{\textbf{Analizy} → \textbf{Regresja} → \textbf{Korelacja}}

\begin{itemize}
\item[$\triangleright$] \textbf{Zmienne}: 2 zmienne = jedna para; 3 i więcej = macierz
\end{itemize}
\begin{enumerate}
\item Współczynnik: r Pearsona, ρ Spearmana (monotoniczna, rangi) albo τ-b Kendalla.
\item Wykres rozrzutu — włączony.
\item Zaawansowane → Hipoteza: Korelacja \textgreater{} 0 / \textless{} 0 (test jednostronny); Przedział ufności.
\item Macierz: Zaawansowane → N w macierzy, Oznacz istotne gwiazdkami.
\end{enumerate}
\textit{Wynik:} Para: jeden wiersz (r, p, N); macierz: dolny trójkąt z r i p w komórkach.

{\color{upwrsoft}\small\textit{Uwaga:} Korelacja nie mówi o przyczynie — sprawdź, czy związek nie wynika z trzeciej zmiennej.}
\end{jbox}

\textbf{Także na tym wykładzie:} Rozkłady t, χ², F, wykładniczy
(s.\textasciitilde{}\pageref{w:rozklady-ciagle}); Przedział ufności dla
różnicy średnich (dwie grupy)
(s.\textasciitilde{}\pageref{w:ci-roznica-srednich}); Przedział ufności
dla średniej różnic (pomiary sparowane)
(s.\textasciitilde{}\pageref{w:ci-srednia-roznic}); Przedział ufności
dla proporcji (s.\textasciitilde{}\pageref{w:ci-proporcja}).

\needspace{8\baselineskip}

\subsection{Wykład 05 --- Założenia
testów}\label{wykux142ad-05-zaux142oux17cenia-testuxf3w}

\begin{jbox}{Sprawdzenie normalności (Shapiro-Wilk, wykres Q-Q)}\label{w:normalnosc}
\jmenu{\textbf{Analizy} → \textbf{Eksploracja} → \textbf{Zmienne ilościowe}}

\begin{itemize}
\item[$\triangleright$] \textbf{Zmienne}: zmienna ilościowa
\item[$\triangleright$] \textbf{Podziel według}: (opcjonalnie) grupy — normalność sprawdza się w każdej grupie
\end{itemize}
\begin{enumerate}
\item Założenia → Normalność: Shapiro-Wilk, Wykres Q-Q.
\end{enumerate}
\textit{Wynik:} Tabela testu (W, p) i wykres Q-Q: punkty blisko prostej = rozkład bliski normalnemu.

{\color{upwrsoft}\small\textit{Uwaga:} W testach t to samo jest w sekcji Sprawdzenie założeń; w ANOVA i regresji sprawdza się reszty (Założenia → Normalność reszt).}
\end{jbox}

\begin{jbox}{Testy nieparametryczne (rangowe)}\label{w:nieparametryczne}
\jmenu{\textbf{Analizy} → \textbf{Testy t / ANOVA}}

\begin{enumerate}
\item Jedna próba, Próby sparowane: Testy → Wilcoxona (rangowanych znaków).
\item Dwie grupy: Testy → Manna-Whitneya U.
\item ANOVA (jeden czynnik): Test → Nieparametrycznie (rangi) = Kruskal-Wallis z porównaniami Dunna.
\item ANOVA powtórzonych pomiarów (jeden czynnik): Nieparametrycznie (rangi) = Friedman.
\end{enumerate}
\textit{Wynik:} Dodatkowy wiersz albo tabela testu rangowego obok wersji parametrycznej.

\end{jbox}

\begin{jbox}{Jednorodność wariancji (Levene, Bartlett)}\label{w:wariancje}
\jmenu{\textbf{Analizy} → \textbf{Testy t / ANOVA}}

\begin{enumerate}
\item Testy t → Dwie grupy → Sprawdzenie założeń: Jednorodność wariancji (Levene).
\item ANOVA → Założenia: Jednorodność wariancji (test Levene’a i Bartletta).
\item Przy różnych wariancjach: t Welcha albo ANOVA → Nierówne wariancje (Welch).
\end{enumerate}
\textit{Wynik:} Tabela testu z F i p.

\end{jbox}

\begin{jbox}{Diagnostyka regresji (reszty, VIF)}\label{w:jakosc-modelu}
\jmenu{\textbf{Analizy} → \textbf{Regresja} → \textbf{Regresja liniowa}}

\begin{enumerate}
\item Założenia → Reszty: Normalność reszt (Shapiro-Wilk), Wykres Q-Q reszt, Reszty vs dopasowane (lejek = niejednorodna wariancja).
\item Założenia → Predyktory i obserwacje: Autokorelacja (Durbin-Watson), Współliniowość (VIF), Odległość Cooka.
\end{enumerate}
\textit{Wynik:} Tabele testów i wykresy reszt pod modelem.

{\color{upwrsoft}\small\textit{Uwaga:} Testu Breuscha-Pagana w jUPWR nie ma — niejednorodność wariancji ocenia się z wykresu Reszty vs dopasowane.}
\end{jbox}

\textbf{Także na tym wykładzie:} Histogram
(s.\textasciitilde{}\pageref{w:histogram}); Test χ² niezależności
(tabela kontyngencji)
(s.\textasciitilde{}\pageref{w:chi2-niezaleznosc}); Mała tabela ---
liczebności oczekiwane i test Fishera
(s.\textasciitilde{}\pageref{w:fisher}); Test t dla dwóch grup
niezależnych (s.\textasciitilde{}\pageref{w:t-dwie-grupy}); ANOVA
jednoczynnikowa i porównania post-hoc
(s.\textasciitilde{}\pageref{w:anova}); Korelacja (Pearson, Spearman) i
macierz korelacji (s.\textasciitilde{}\pageref{w:korelacja}).

\needspace{8\baselineskip}

\subsection{Wykład 06 --- Regresja}\label{wykux142ad-06-regresja}

\begin{jbox}{Przedział ufności dla korelacji i nachylenia regresji}\label{w:ci-korelacja-regresja}
\jmenu{\textbf{Analizy} → \textbf{Przedziały ufności} → \textbf{Korelacja i regresja} → \textbf{Korelacja / Regresja liniowa}}

\begin{itemize}
\item[$\triangleright$] \textbf{Zmienna 1, Zmienna 2}: dwie zmienne ilościowe (Korelacja)
\item[$\triangleright$] \textbf{Zmienna zależna (Y), Predyktor (X)}: (Regresja liniowa)
\end{itemize}
\begin{enumerate}
\item Korelacja: Współczynnik r Pearsona albo ρ Spearmana; Metoda: Transformacja Fishera z.
\item Regresja liniowa: Metoda t-Studenta (lm).
\end{enumerate}
\textit{Wynik:} Tabela z przedziałem dla r albo dla wyrazu wolnego i nachylenia; wykres z pasmem.

\end{jbox}

\begin{jbox}{Regresja liniowa prosta}\label{w:regresja-prosta}
\jmenu{\textbf{Analizy} → \textbf{Regresja} → \textbf{Regresja liniowa}}

\begin{itemize}
\item[$\triangleright$] \textbf{Zmienna zależna}: Y — zmienna ilościowa
\item[$\triangleright$] \textbf{Predyktory ilościowe}: X
\end{itemize}
\begin{enumerate}
\item Wykres rozrzutu z prostą (jeden predyktor) — włączony.
\item Zaawansowane → Zapis do arkusza: Wartości przewidywane do arkusza.
\end{enumerate}
\textit{Wynik:} Dopasowanie modelu: R², R² skoryg., F, p, RMSE. Współczynniki: b (nachylenie), SE, t, p, przedział ufności.

\end{jbox}

\begin{jbox}{Regresja wieloraka (także z predyktorem jakościowym)}\label{w:regresja-wieloraka}
\jmenu{\textbf{Analizy} → \textbf{Regresja} → \textbf{Regresja liniowa}}

\begin{itemize}
\item[$\triangleright$] \textbf{Zmienna zależna}: Y
\item[$\triangleright$] \textbf{Predyktory ilościowe}: zmienne ilościowe
\item[$\triangleright$] \textbf{Predyktory jakościowe}: zmienne jakościowe (np. płeć, gatunek)
\end{itemize}
\begin{enumerate}
\item Poziomy odniesienia predyktorów jakościowych: wybierz poziom, z którym porównujesz pozostałe.
\item Zaawansowane → Model: Współczynniki standaryzowane β, AIC i BIC.
\item Założenia: Współliniowość (VIF).
\end{enumerate}
\textit{Wynik:} Współczynniki: każdy b = zmiana Y przy wzroście predyktora o 1, gdy pozostałe są stałe; R² skoryg. do porównania modeli.

\end{jbox}

\begin{jbox}{Regresja logistyczna}\label{w:logistyczna}
\jmenu{\textbf{Analizy} → \textbf{Regresja} → \textbf{Regresja logistyczna} → \textbf{Regresja logistyczna}}

\begin{itemize}
\item[$\triangleright$] \textbf{Zmienna zależna (2 poziomy)}: zmienna dwukategorialna (np. zdał / nie zdał)
\item[$\triangleright$] \textbf{Predyktory ilościowe}: zmienne ilościowe
\item[$\triangleright$] \textbf{Predyktory jakościowe}: zmienne jakościowe
\end{itemize}
\begin{enumerate}
\item Kategoria „zdarzenie”: poziom, którego szansę modelujesz (domyślnie drugi alfabetycznie).
\item Zaawansowane → Klasyfikacja: Próg klasyfikacji = 0.5, Krzywa ROC.
\end{enumerate}
\textit{Wynik:} Współczynniki: b, p i iloraz szans (OR) z przedziałem; tabela klasyfikacji z trafnością, czułością i swoistością.

{\color{upwrsoft}\small\textit{Uwaga:} Zmienną 0/1 z ilościowej utworzysz przez Dane → Oblicz, np. read \textgreater{}= VMED(read).}
\end{jbox}

\textbf{Także na tym wykładzie:} Nowa zmienna z formuły
(s.\textasciitilde{}\pageref{w:oblicz-zmienna}); Diagnostyka regresji
(reszty, VIF) (s.\textasciitilde{}\pageref{w:jakosc-modelu}).

\needspace{8\baselineskip}

\subsection{Wykład 07 --- Dobre dane}\label{wykux142ad-07-dobre-dane}

\begin{jbox}{Wykluczenie obserwacji (filtr)}\label{w:filtry}
\jmenu{\textbf{Dane} → \textbf{Filtry}}

\begin{enumerate}
\item Wpisz warunek, który mają SPEŁNIAĆ pozostawione wiersze, np. wiek \textless{} 100 albo gatunek != 'Adelie'.
\item Kolejny warunek: Dodaj nowy filtr. Wyłączenie bez kasowania: Filtr jest aktywny.
\end{enumerate}
\textit{Wynik:} Odfiltrowane wiersze są wyszarzone i znikają ze wszystkich analiz.

\end{jbox}

\begin{jbox}{Braki danych i wartości skrajne}\label{w:braki-odstajace}
\jmenu{\textbf{Analizy} → \textbf{Eksploracja} → \textbf{Zmienne ilościowe}}

\begin{itemize}
\item[$\triangleright$] \textbf{Zmienne}: wszystkie zmienne ilościowe
\end{itemize}
\begin{enumerate}
\item Braki: wiersz Braki w tabeli (domyślnie).
\item Zaawansowane → Wartości skrajne: pięć najmniejszych i największych wartości z numerami wierszy.
\item Wykres: Pudełkowy z Punktami.
\end{enumerate}
\textit{Wynik:} Tabela wartości skrajnych wskazuje wiersze do sprawdzenia w arkuszu.

{\color{upwrsoft}\small\textit{Uwaga:} Dla zmiennych jakościowych braki pokazuje Eksploracja → Zmienne jakościowe → Podsumowanie.}
\end{jbox}

\textbf{Także na tym wykładzie:} Otwarcie pliku z danymi
(s.\textasciitilde{}\pageref{w:otworz-plik}); Typ zmiennej (poziom
pomiaru) (s.\textasciitilde{}\pageref{w:poziom-pomiaru}); Nowa zmienna z
formuły (s.\textasciitilde{}\pageref{w:oblicz-zmienna}); Średnia,
mediana, kwartyle, odchylenie standardowe
(s.\textasciitilde{}\pageref{w:statystyki-opisowe}); Wykres pudełkowy
(s.\textasciitilde{}\pageref{w:pudelkowy}).

\needspace{8\baselineskip}

\subsection{Wykład 08 --- Case study --- szkoły w
Kalifornii}\label{wykux142ad-08-case-study-szkoux142y-w-kalifornii}

\textbf{Także na tym wykładzie:} Średnia, mediana, kwartyle, odchylenie
standardowe (s.\textasciitilde{}\pageref{w:statystyki-opisowe});
Histogram (s.\textasciitilde{}\pageref{w:histogram}); Wykres pudełkowy
(s.\textasciitilde{}\pageref{w:pudelkowy}); Korelacja (Pearson,
Spearman) i macierz korelacji
(s.\textasciitilde{}\pageref{w:korelacja}); Regresja liniowa prosta
(s.\textasciitilde{}\pageref{w:regresja-prosta}); Regresja wieloraka
(także z predyktorem jakościowym)
(s.\textasciitilde{}\pageref{w:regresja-wieloraka}).

\needspace{8\baselineskip}

\subsection{Wykład 09 --- Projekt
badawczy}\label{wykux142ad-09-projekt-badawczy}

\begin{jbox}{Opis metod do raportu}\label{w:opis-metod}
\jmenu{\textbf{(dowolna analiza jUPWR)} → \textbf{Opis metod}}

\begin{enumerate}
\item Zaznacz „Opis zastosowanych metod” — na końcu wyników pojawi się opis: dane, test, parametry, wielkość efektu.
\end{enumerate}
\textit{Wynik:} Akapit pod tabelami analizy; opisuje metody, nie wyniki.

\end{jbox}

\textbf{Także na tym wykładzie:} Otwarcie pliku z danymi
(s.\textasciitilde{}\pageref{w:otworz-plik}); Nowa zmienna z formuły
(s.\textasciitilde{}\pageref{w:oblicz-zmienna}); Wykluczenie obserwacji
(filtr) (s.\textasciitilde{}\pageref{w:filtry}); Średnia, mediana,
kwartyle, odchylenie standardowe
(s.\textasciitilde{}\pageref{w:statystyki-opisowe}); Statystyki opisowe
osobno w grupach (s.\textasciitilde{}\pageref{w:porownanie-grup-opis});
Test t dla dwóch grup niezależnych
(s.\textasciitilde{}\pageref{w:t-dwie-grupy}); Testy nieparametryczne
(rangowe) (s.\textasciitilde{}\pageref{w:nieparametryczne}); Korelacja
(Pearson, Spearman) i macierz korelacji
(s.\textasciitilde{}\pageref{w:korelacja}); Regresja wieloraka (także z
predyktorem jakościowym)
(s.\textasciitilde{}\pageref{w:regresja-wieloraka}).

\newpage

\section{Część B --- pytanie →
analiza}\label{czux119ux15bux107-b-pytanie-analiza}

\needspace{6\baselineskip}

\subsection{Przygotowanie danych}\label{przygotowanie-danych}

\begin{longtable*}{@{}p{6.4cm}p{9.2cm}r@{}}
\toprule
\textbf{Pytanie} & \textbf{Gdzie w jUPWR} & \textbf{s.} \\
\midrule
\endhead
Jak wczytać dane (CSV, Excel, .omv)? & \jmenu{\textbf{{\symfont ☰} (menu pliku)} → \textbf{Otwórz} → \textbf{Przeglądaj}} & \pageref{w:otworz-plik} \\[2pt]
Jak ustawić, czy zmienna jest jakościowa czy ilościowa? & \jmenu{\textbf{Dane} → \textbf{Edytuj}} & \pageref{w:poziom-pomiaru} \\[2pt]
Jak utworzyć zmienną, np. „czy english \textgreater{} 20”, średnią dwóch kolumn albo z-score? & \jmenu{\textbf{Dane} → \textbf{Oblicz}} & \pageref{w:oblicz-zmienna} \\[2pt]
Jak pominąć część obserwacji, np. błędne wartości albo jedną grupę? & \jmenu{\textbf{Dane} → \textbf{Filtry}} & \pageref{w:filtry} \\[2pt]
Jak opisać w raporcie, co zostało policzone? & \jmenu{\textbf{(dowolna analiza jUPWR)} → \textbf{Opis metod}} & \pageref{w:opis-metod} \\[2pt]
Czy w danych są braki albo podejrzane wartości? & \jmenu{\textbf{Analizy} → \textbf{Eksploracja} → \textbf{Zmienne ilościowe}} & \pageref{w:braki-odstajace} \\[2pt]
\bottomrule
\end{longtable*}

\needspace{6\baselineskip}

\subsection{Opis danych}\label{opis-danych}

\begin{longtable*}{@{}p{6.4cm}p{9.2cm}r@{}}
\toprule
\textbf{Pytanie} & \textbf{Gdzie w jUPWR} & \textbf{s.} \\
\midrule
\endhead
Ile jest obserwacji w każdej kategorii? & \jmenu{\textbf{Analizy} → \textbf{Eksploracja} → \textbf{Zmienne jakościowe}} & \pageref{w:tabela-licznosci} \\[2pt]
Jak opisać zmienną ilościową (położenie i rozrzut)? & \jmenu{\textbf{Analizy} → \textbf{Eksploracja} → \textbf{Zmienne ilościowe}} & \pageref{w:statystyki-opisowe} \\[2pt]
Jak porównać opis zmiennej ilościowej między grupami? & \jmenu{\textbf{Analizy} → \textbf{Eksploracja} → \textbf{Zmienne ilościowe}} & \pageref{w:porownanie-grup-opis} \\[2pt]
Jaki kształt ma rozkład zmiennej? & \jmenu{\textbf{Analizy} → \textbf{Eksploracja} → \textbf{Zmienne ilościowe}} & \pageref{w:histogram} \\[2pt]
Gdzie jest mediana, kwartyle i wartości odstające? & \jmenu{\textbf{Analizy} → \textbf{Eksploracja} → \textbf{Zmienne ilościowe}} & \pageref{w:pudelkowy} \\[2pt]
\bottomrule
\end{longtable*}

\needspace{6\baselineskip}

\subsection{Rozkłady}\label{rozkux142ady}

\begin{longtable*}{@{}p{6.4cm}p{9.2cm}r@{}}
\toprule
\textbf{Pytanie} & \textbf{Gdzie w jUPWR} & \textbf{s.} \\
\midrule
\endhead
Ile wynosi P(X \textless{} a) albo P(a \textless{} X \textless{} b) dla rozkładu normalnego? & \jmenu{\textbf{Analizy} → \textbf{Rozkłady} → \textbf{Rozkład normalny}} & \pageref{w:rozklad-normalny} \\[2pt]
Ile wynosi P(X = k) albo P(X ≤ k) w rozkładzie dyskretnym? & \jmenu{\textbf{Analizy} → \textbf{Rozkłady} → \textbf{Rozkład dwumianowy / Poissona / geometryczny}} & \pageref{w:rozklady-dyskretne} \\[2pt]
Czy rozkład zmiennej jest zbliżony do normalnego? & \jmenu{\textbf{Analizy} → \textbf{Eksploracja} → \textbf{Zmienne ilościowe}} & \pageref{w:normalnosc} \\[2pt]
\bottomrule
\end{longtable*}

\needspace{6\baselineskip}

\subsection{Oszacowanie z niepewnością (przedziały
ufności)}\label{oszacowanie-z-niepewnoux15bciux105-przedziaux142y-ufnoux15bci}

\begin{longtable*}{@{}p{6.4cm}p{9.2cm}r@{}}
\toprule
\textbf{Pytanie} & \textbf{Gdzie w jUPWR} & \textbf{s.} \\
\midrule
\endhead
Jaka jest średnia w populacji (z niepewnością)? & \jmenu{\textbf{Analizy} → \textbf{Przedziały ufności} → \textbf{Średnie} → \textbf{Średnia}} & \pageref{w:ci-srednia} \\[2pt]
O ile różnią się średnie dwóch grup (z niepewnością)? & \jmenu{\textbf{Analizy} → \textbf{Przedziały ufności} → \textbf{Średnie} → \textbf{Różnica średnich}} & \pageref{w:ci-roznica-srednich} \\[2pt]
Jaki jest udział kategorii w populacji (z niepewnością)? & \jmenu{\textbf{Analizy} → \textbf{Przedziały ufności} → \textbf{Proporcje} → \textbf{Proporcja}} & \pageref{w:ci-proporcja} \\[2pt]
\bottomrule
\end{longtable*}

\needspace{6\baselineskip}

\subsection{Testy --- porównania}\label{testy-poruxf3wnania}

\begin{longtable*}{@{}p{6.4cm}p{9.2cm}r@{}}
\toprule
\textbf{Pytanie} & \textbf{Gdzie w jUPWR} & \textbf{s.} \\
\midrule
\endhead
Czy średnia różni się od zadanej wartości (np. 650)? & \jmenu{\textbf{Analizy} → \textbf{Testy t} → \textbf{Jedna próba}} & \pageref{w:t-jedna-proba} \\[2pt]
Czy udział kategorii różni się od zadanej wartości (np. 50\%)? & \jmenu{\textbf{Analizy} → \textbf{Częstości} → \textbf{Jedna zmienna}} & \pageref{w:test-proporcji} \\[2pt]
Czy rozkład kategorii zgadza się z oczekiwanym (np. równe udziały)? & \jmenu{\textbf{Analizy} → \textbf{Częstości} → \textbf{Jedna zmienna}} & \pageref{w:chi2-zgodnosc} \\[2pt]
Czy dwie zmienne jakościowe są powiązane? & \jmenu{\textbf{Analizy} → \textbf{Częstości} → \textbf{Tabela kontyngencji}} & \pageref{w:chi2-niezaleznosc} \\[2pt]
Co zrobić, gdy liczebności w tabeli są małe? & \jmenu{\textbf{Analizy} → \textbf{Częstości} → \textbf{Tabela kontyngencji}} & \pageref{w:fisher} \\[2pt]
Czy średnie w dwóch grupach się różnią? & \jmenu{\textbf{Analizy} → \textbf{Testy t} → \textbf{Dwie grupy}} & \pageref{w:t-dwie-grupy} \\[2pt]
Czy wynik zmienił się między dwoma pomiarami tych samych osób? & \jmenu{\textbf{Analizy} → \textbf{Testy t} → \textbf{Próby sparowane}} & \pageref{w:t-sparowane} \\[2pt]
Czy średnie w trzech lub więcej grupach się różnią — i które? & \jmenu{\textbf{Analizy} → \textbf{ANOVA} → \textbf{ANOVA}} & \pageref{w:anova} \\[2pt]
Czego użyć, gdy założenia testu t / ANOVA nie są spełnione? & \jmenu{\textbf{Analizy} → \textbf{Testy t / ANOVA}} & \pageref{w:nieparametryczne} \\[2pt]
Czy wariancje w grupach są podobne? & \jmenu{\textbf{Analizy} → \textbf{Testy t / ANOVA}} & \pageref{w:wariancje} \\[2pt]
Jak duży jest efekt (a nie tylko: czy istotny)? & \jmenu{\textbf{Analizy} → \textbf{(zależnie od testu)}} & \pageref{w:sila-efektu} \\[2pt]
\bottomrule
\end{longtable*}

\needspace{6\baselineskip}

\subsection{Zależności między
zmiennymi}\label{zaleux17cnoux15bci-miux119dzy-zmiennymi}

\begin{longtable*}{@{}p{6.4cm}p{9.2cm}r@{}}
\toprule
\textbf{Pytanie} & \textbf{Gdzie w jUPWR} & \textbf{s.} \\
\midrule
\endhead
Czy dwie zmienne ilościowe są powiązane — jak silnie i w którą stronę? & \jmenu{\textbf{Analizy} → \textbf{Regresja} → \textbf{Korelacja}} & \pageref{w:korelacja} \\[2pt]
Jak zmienia się Y, gdy X rośnie o jednostkę? Jak przewidzieć Y? & \jmenu{\textbf{Analizy} → \textbf{Regresja} → \textbf{Regresja liniowa}} & \pageref{w:regresja-prosta} \\[2pt]
Czy model regresji jest poprawny? & \jmenu{\textbf{Analizy} → \textbf{Regresja} → \textbf{Regresja liniowa}} & \pageref{w:jakosc-modelu} \\[2pt]
Jak kilka zmiennych naraz wpływa na Y? & \jmenu{\textbf{Analizy} → \textbf{Regresja} → \textbf{Regresja liniowa}} & \pageref{w:regresja-wieloraka} \\[2pt]
Od czego zależy szansa zdarzenia (zmienna tak/nie)? & \jmenu{\textbf{Analizy} → \textbf{Regresja} → \textbf{Regresja logistyczna} → \textbf{Regresja logistyczna}} & \pageref{w:logistyczna} \\[2pt]
\bottomrule
\end{longtable*}




\end{document}
